\Needspace{4\baselineskip}
\begin{Example}{波动方程}
\label{ex:wave-fundamental-solution}
对二阶方程
\begin{equation}\label{eq:second-order-time-evolution-fundamental-equation}
u_{tt}=L(D)u
\end{equation}
将基本解 $G$ 定义为二次连续可微的函数 $[0,\infty)\to\mathcal{D}^{\prime}(\mathbb{R}^m)$, 使 $G$ 满足 \eqref{eq:second-order-time-evolution-fundamental-equation} 及初始条件 $G(\bm{x},0)=0$, $G_t(\bm{x},0)=\delta(\bm{x})$. 那么非齐次问题
\begin{equation}\label{eq:second-order-evolution-inhomogeneous-initial-problem}
u_{tt}=L(D)u+f(\bm{x},t),\quad u(\bm{x},0)=u_0(\bm{x}),\quad u_t(\bm{x},0)=u_1(\bm{x})
\end{equation}
的解表示为
\begin{equation}\label{eq:second-order-evolution-fundamental-solution-formula}
\begin{aligned}
u(\bm{x},t)&=\int_{\mathbb{R}^m}G_t(\bm{x}-\bm{y},t)u_0(\bm{y})\,\mathrm{d}\bm{y}
+\int_{\mathbb{R}^m}G(\bm{x}-\bm{y},t)u_1(\bm{y})\,\mathrm{d}\bm{y}\\
&\quad+\int_0^t\int_{\mathbb{R}^m}G(\bm{x}-\bm{y},t-s)f(\bm{y},s)\,\mathrm{d}\bm{y}\,\mathrm{d}s.
\end{aligned}
\end{equation}
对一维空间中的波动方程,
\begin{equation}\label{eq:wave-one-dimensional-fundamental-solution}
G(x,t)=\begin{cases}1/2&|x|<t\\0&|x|\geq t\end{cases}
\end{equation}
是一个基本解. 事实上, 显然有 $G(x,0)=0$, 并且由表示式 $G(x,t)=(H(x+t)-H(x-t))/2$ 可得, $G$ 满足波动方程, 且 $G_t(x,t)=(\delta(x+t)+\delta(x-t))/2$, 当 $t=0$ 时它等于 $\delta(x)$. 多维波动方程的基本解将在下一节讨论. 提请注意, Laplace 方程和热方程的基本解 (除原点处的奇性外) 是 $C^\infty$ 函数, 但波动方程的基本解不是. 这对解的正则性有重要影响.
\end{Example}

\paragraph*{Problems}

\begin{enumerate}[label=\thesection.\arabic*.,ref=\thesection.\arabic*,start=33,series=book-problems]

\item\label{ex:distribution-smooth-test-parameter-map} 设 $\phi(\bm{x},\bm{y})\in\mathcal{D}(\mathbb{R}^{p+q})$, $g\in\mathcal{D}^{\prime}(\mathbb{R}^q)$. 证明 \[\psi(\bm{x}):=(g(\bm{y}),\phi(\bm{x},\bm{y}))\] 属于 $\mathcal{D}(\mathbb{R}^p)$.

\item\label{ex:tempered-direct-product} 证明在 $\mathcal{S}^{\prime}$ 中也可定义直积.

\item\label{ex:direct-product-support} 设 $F$ 和 $G$ 为 $f$ 和 $g$ 的支集. 证明直积的支集为 $F\times G$.

\item\label{ex:test-function-integral-riemann-sums} 设 $\phi,\psi\in\mathcal{D}(\mathbb{R}^m)$. 证明 $\int_{\mathbb{R}^m}\psi(\bm{y}-\bm{x})\phi(\bm{y})\,\mathrm{d}\bm{y}$ 定义了 $\mathcal{D}(\mathbb{R}^m)$ 的一个元素. 此外, 证明在 $\mathcal{D}(\mathbb{R}^m)$ 中的收敛意义下, 该积分是 Riemann 和的极限.

\item\label{ex:distribution-convolution-continuity} 证明 Lemma~\ref{lem:distribution-convolution-continuity}.

\item\label{ex:test-functions-dense-tempered-distributions} 讨论如何修改 Theorem~\ref{thm:test-functions-dense-distributions} 的证明, 以说明 $\mathcal{D}(\Omega)$ 在 $\mathcal{D}^{\prime}(\Omega)$ 中稠密. 再证明 $\mathcal{D}(\mathbb{R}^m)$ 在 $\mathcal{S}^{\prime}(\mathbb{R}^m)$ 中稠密.

\item\label{ex:direct-product-joint-continuity} 证明直积关于其因子联合连续.

\item\label{ex:biharmonic-fundamental-solution} 求 $\mathbb{R}^m$ 上双调和方程 $\Delta\Delta G=\delta(\bm{x})$ 的一个基本解.

\item\label{ex:inhomogeneous-fundamental-solution-formula} 证明 \eqref{eq:first-order-evolution-fundamental-solution-formula} 确实是 \eqref{eq:first-order-evolution-inhomogeneous-initial-problem} 的解.

\item\label{ex:evolution-fundamental-spacetime-delta} 设 $G$ 为 $u_t=L(D)u$ 的初值问题所对应的基本解. 证明泛函
\begin{equation}\label{eq:evolution-fundamental-solution-spacetime-functional}
F:\phi\mapsto\int_0^\infty\int_{\mathbb{R}^m}G(\bm{x},t)\phi(\bm{x},t)\,\mathrm{d}\bm{x}\,\mathrm{d}t
\end{equation}
定义了 $\mathbb{R}^{m+1}$ 上的分布, 且该分布满足方程 $F_t-L(D)F=\delta(\bm{x},t)$.

\item\label{ex:wave-one-dimensional-duhamel-formula} 将 \eqref{eq:second-order-evolution-fundamental-solution-formula} 专门用于一维波动方程.

\item\label{ex:polynomial-compact-distribution-convolution} 设 $f$ 为紧支集分布, 且 $P$ 为多项式. 证明 $P*f$ 为多项式.

\end{enumerate}
